\section{Related Work}

\subsection{Peptide--MHC Binding and Presentation Prediction}

Peptide--MHC binding prediction is a foundational problem in computational immunology. Established methods such as NetMHCpan \cite{reynisson2020netmhcpan}, MHCflurry \cite{odonnell2020mhcflurry}, and ACME \cite{hu2019acme} estimate binding affinity, binding rank, or peptide--MHC compatibility from peptide sequences and allele information. Pan-allelic predictors improve coverage by representing MHC molecules through pseudo-sequences that summarize the peptide-binding groove or through related allele features. This representation allows information sharing across well-characterized and data-scarce alleles \cite{nielsen2007netmhcpan,reynisson2020netmhcpan}. These models are essential for filtering candidate epitopes, but binding is only one stage of antigen presentation. A peptide with favorable binding characteristics may still be absent from the cell-surface immunopeptidome because its source protein is not expressed or degraded, the peptide is not generated or transported efficiently, or it is not detected experimentally. Binding scores therefore do not directly encode presentation-evidence preferences among tissues that share the same MHC restriction.

\subsection{Predicting Cell-surface Antigen Presentation (Antigen Presentation Prediction)}

Eluted-ligand data enable models to learn a broader presentation signal than binding measurements alone \cite{reynisson2020netmhcpan,odonnell2020mhcflurry}. General presentation predictors combine peptide sequence, MHC context, and, in some cases, antigen-processing or source-protein features to distinguish observed ligands from background peptides. Systems such as BigMHC \cite{albert2023bigmhc} illustrate how large immunopeptidomic collections and modern sequence models can improve ligand presentation prediction and candidate ranking. Their usual target, however, is whether a peptide is likely to be presented for an allele, sample, or pooled presentation setting. This target differs from the relative question studied here: given a particular tissue--MHC context, which of two peptides from the same parent protein has stronger evidence of tissue-specific preference in presentation? Consequently, a high general-presentation score need not imply a tissue-specific preference, and direct comparisons require a benchmark whose supervision reflects that distinction.

\subsection{Tissue- and Sample-Context Immunopeptidomics}

Multi-tissue immunopeptidomic atlases have revealed that ligand repertoires contain both broadly shared peptides and tissue-associated components \cite{marcu2021atlas,schuster2018mouse}. MHCflurry 2.0 incorporates antigen-processing features but does not use tissue identity \cite{odonnell2020mhcflurry}. More recent systems perform multi-allelic deconvolution, as in ImmuneApp \cite{xu2024immuneapp}, or combine peptide, MHC, peptide-flanking sequence, and additional molecular context, as in HLApollo \cite{thrift2024hlapollo}. These approaches demonstrate that presentation is context dependent, but they generally do not define a tissue--MHC combination as the supervised prediction unit or match every recorded positive and pseudo-negative on both MHC restriction and parent protein. Our formulation instead isolates a relative-ranking problem designed to focus on tissue-associated antigen-processing preference.

\subsection{Sharing Information across Tissue--MHC Combinations}

Some tissue--MHC combinations contain many observations, whereas others contain relatively few. Learning a completely separate model for every combination would waste information shared across tissues and alleles. Conversely, pooling all combinations into one model could blur allele-specific binding motifs and tissue-associated differences. Joint learning provides a middle ground: related combinations share a peptide representation while retaining separate prediction functions \cite{caruana1997multitask,ruder2017overview}. More structured methods allow different groups to use shared and specialized model components \cite{ma2018mmoe,tang2020ple}. TissuePMHC adopts a biologically structured sharing scheme: one model view learns patterns across all human combinations, while a second shares information only among tissues with the same HLA allele.

\subsection{Relation to This Work}

The principal distinction from previous work is the prediction target. General predictors estimate peptide binding to or presentation by an \plainMHC{}, whereas TissuePMHC ranks peptides matched by MHC restriction, parent protein, and total positive tissue-occurrence count within a specified tissue--MHC context. The human and mouse datasets define this comparison identically, although each species uses a separately trained model. The current occurrence-matched experiments use a fixed internal test set and allow a peptide to recur in another fitting combination. MHCflurry and NetMHCpan are therefore used as tissue-independent peptide--MHC controls rather than as tissue-aware competitors.

These distinctions are stated in prose because the published methods use different targets, negative constructions, datasets, and evaluation protocols. NetMHCpan and MHCflurry provide tissue-independent peptide--MHC controls in the present fixed test; BigMHC, ImmuneApp, and HLApollo remain related presentation systems rather than directly comparable tissue-aware baselines. The controlled architecture surveys below therefore contain only systems evaluated on the same occurrence-matched rows, task inventory, and fixed split.

Model attribution addresses a separate question from predictive benchmarking: which input features support a fitted model's output. SHAP provides an additive feature-attribution framework, while integrated gradients attributes a differentiable output relative to a baseline along an input path \cite{lundberg2017shap,sundararajan2017integrated}. Expected gradients combines integration with a distribution of reference examples and provides a practical SHAP-style approximation for differentiable networks \cite{erion2021expected}. In silico mutagenesis (ISM) instead perturbs sequence symbols and measures the resulting prediction change, providing a direct model-response check that does not depend on a gradient path or attribution background \cite{alipanahi2015deepbind}. We use attribution and ISM post hoc, after model selection and fixed-test prediction, to describe and cross-check position- and residue-level dependence of the fitted Human branches; neither analysis is used to select a model or infer a causal antigen-processing pathway.
